% System Description — COUNTED in 15-page limit
% Architecture of final system. Flow charts, schematics, images.
% Merges: system_architecture details, hardware specs, CV pipeline, target localisation,
%         ground station, simulation

\section{System Description}\label{sec:sysdesc}

The final system integrates seven subsystems into an autonomous SAR platform: a hardware airframe with companion computer (\S\ref{sec:hw}), eleven Python modules (\S\ref{sec:sw-arch}), a YOLOv8n vision pipeline at 4.8\,FPS (\S\ref{sec:cv}), a boustrophedon search pattern with formal coverage guarantee (\S\ref{sec:search-pattern}), a 20-state mission FSM with six safety mechanisms (\S\ref{sec:state-machine}), GPS target localisation at CEP$_{50}$ = 2.3\,m (\S\ref{sec:localisation}), and a browser-based ground station at \SI{300}{\milli\second} video latency (\S\ref{sec:groundstation}). The design philosophy is \emph{independent modules composed at runtime}: vision, planning, geofencing, and navigation have zero mutual dependencies, enabling isolated testing and single-file swaps with no integration risk. Figure~\ref{fig:mission_overview} shows the complete mission plan overlaid on the Fenswood Farm site.

\begin{figure}[ht]
\centering
\includegraphics[width=\linewidth]{mission_overview}
\caption{Mission overview on the Fenswood Farm site. The lawnmower search pattern covers the designated survey polygon at \SI{35}{\metre} altitude, with the take-off point, return-to-launch path, and SSSI no-fly zone boundary shown. The flight area boundary constrains all autonomous waypoints.}
\label{fig:mission_overview}
\end{figure}

% ═══════════════════════════════════════════════════════════════
\subsection{Hardware Platform}\label{sec:hw}

Table~\ref{tab:hw-bom} summarises the bill of materials for the Hexsoon EDU-450 airframe, CubePilot Orange+ flight controller (ArduPilot Copter 4.5), Raspberry Pi~5 companion computer, and IMX296 global shutter camera. All hardware was provided by the university as part of the AENGM0074 equipment package; design effort therefore focused on system integration, software architecture, and calibration rather than component selection (see Appendix~\ref{sec:cv:inference} for the edge inference architecture rationale).

\begin{table}[ht]
  \centering\compactTable
  \caption{System bill of materials and approximate costs.}\label{tab:hw-bom}
  \begin{tabularx}{\linewidth}{@{}l X r@{}}
    \toprule
    \textbf{Component} & \textbf{Specification} & \textbf{Cost (\pounds)} \\
    \midrule
    Cube Orange+ FC   & STM32H757, triple-redundant IMU, ArduPilot 4.5 & 180 \\
    Here3 GPS          & u-blox M8P, CAN interface, 5\,Hz updates & 70 \\
    Raspberry Pi 5     & BCM2712 quad-core A76, \SI{8}{\giga\byte} RAM & 60 \\
    IMX296 GS Camera   & Sony IMX296, $1456\!\times\!1088$, global shutter, CSI-2 & 45 \\
    Frame + motors + ESCs & EDU-450, 2212/920\,KV brushless, 30\,A BLHeli\_S & 100 \\
    Battery + RC + misc & 4S \SI{5200}{\milli\ampere\hour} LiPo, FrSky Twin X14, wiring & 110 \\
    \midrule
    \textbf{Total}     & & \textbf{$\sim$565} \\
    \bottomrule
  \end{tabularx}
\end{table}

The Pi communicates with the Cube via UART at 921\,600\,baud through a MAVProxy bridge, which republishes MAVLink traffic over UDP. A second output on TCP port~5762 allows Mission Planner on the ground station laptop to monitor telemetry via Wi-Fi (see Appendix~\ref{sec:streaming-arch} for the full communications architecture). The camera captures at $1456 \times 1088$ via CSI-2; lens distortion is corrected using precomputed remap matrices (calibrated RMS = 0.399\,px, $+$1.5\,ms overhead; methodology in Appendix~\ref{sec:calibration}). The calibrated focal length is \SI{5.46}{\milli\metre}, yielding a horizontal field of view (HFOV) of ${\approx}\,49.3^{\circ}$ (Pi IMX296 camera).

% ═══════════════════════════════════════════════════════════════
\subsection{Software Architecture}\label{sec:sw-arch}

The system comprises eleven Python modules totalling approximately 4\,400 lines of production code, organised in four layers (Figure~\ref{fig:architecture}). An additional 74 test and calibration scripts (Appendix~\ref{sec:test-scripts}), including 127 unit tests, validate each module in isolation before integration:

\begin{figure}[ht]
\centering
\includegraphics[width=0.85\textwidth]{architecture}
\caption{Module dependency graph. Independent modules have no cross-dependencies; the orchestrator composes them at runtime.}
\label{fig:architecture}
\end{figure}

\begin{figure}[ht]
\centering
\includegraphics[width=0.75\textwidth]{pi_system}
\caption{Raspberry Pi~5 companion computer system: hardware interfaces and data flow between the CSI camera, TFLite inference engine, MAVProxy bridge, and ground station.}
\label{fig:pi_system}
\end{figure}

\begin{itemize}
    \item \textbf{Configuration layer:} \texttt{config.py} centralises all tuneable parameters (Appendix~\ref{app:config}); \texttt{states.py} defines the 20-state enumeration.
    \item \textbf{Independent modules:} \texttt{vision.py} (camera + AI detection), \texttt{planning.py} (lawnmower pattern), \texttt{gps\_utils.py} (coordinate transforms), \texttt{geofence.py} (NFZ enforcement). These have no mutual dependencies and can be tested in isolation.
    \item \textbf{Domain modules:} \texttt{navigation.py} wraps MAVLink commands (Table~\ref{tab:mavlink-commands}); \texttt{stream\_server.py} serves the ground station.
    \item \textbf{Orchestration layer:} \texttt{state\_machine.py} implements all state handlers as a mixin class; \texttt{main.py} composes everything and runs the main loop.
\end{itemize}

\begin{table}[ht]
\centering\small
\begin{tabular}{lll}
\toprule
\textbf{Method} & \textbf{MAVLink Message} & \textbf{Purpose} \\
\midrule
\texttt{arm()} & \texttt{MAV\_CMD\_COMPONENT\_ARM\_DISARM} & Arm/disarm motors \\
\texttt{takeoff(alt)} & \texttt{MAV\_CMD\_NAV\_TAKEOFF} & Climb to altitude \\
\texttt{send\_global\_target()} & \texttt{SET\_POSITION\_TARGET\_GLOBAL\_INT} & Fly to GPS waypoint \\
\texttt{send\_velocity()} & \texttt{SET\_POSITION\_TARGET\_LOCAL\_NED} & Velocity command \\
\texttt{set\_mode()} & \texttt{SET\_MODE} & Change flight mode \\
\texttt{land()} & \texttt{MAV\_CMD\_NAV\_LAND} & Land at position \\
\bottomrule
\end{tabular}
\caption{MAVLink commands used by the navigation module.}
\label{tab:mavlink-commands}
\end{table}

The data pipeline traverses five stages per iteration: (1)~capture a frame and apply lens undistortion; (2)~resize to $640\times640$, normalise, and run YOLOv8n inference; (3)~convert the pixel detection to a GPS estimate; (4)~dispatch the current state handler; and (5)~enforce geofence constraints. Geofencing always runs \emph{after} the state handler, so waypoint commands can be overridden by repulsive velocities.

\paragraph{Geofence enforcement.}
The geofence module implements five independent protection layers (Figure~\ref{fig:geofence_diagram}). The outermost layer is the \textbf{Flight Area boundary}: a four-vertex polygon checked via \texttt{cv2.pointPolygonTest} every main-loop iteration (${\sim}$20--50\,Hz); any excursion triggers an immediate emergency RTL with no buffer or warning zone. The remaining four layers protect the \textbf{SSSI no-fly zone}: (1)~plan-time waypoint filtering (30\,m buffer); (2)~a speed scalar field that linearly decelerates the drone from 3.0\,m/s at 20\,m to a full stop at 2\,m from the SSSI boundary; (3)~a repulsive potential field that pushes the aircraft away at 5.0\,m/s when within ${\sim}$3\,m of the boundary in MANUAL mode; and (4)~a hard cutoff that forces MANUAL mode if the drone enters the SSSI polygon. The module computes the repulsive offset using \texttt{cv2.pointPolygonTest} for signed distance and applies it as a velocity correction after the state handler's waypoint command. This layered design ensures no state can inadvertently fly into restricted airspace (see Appendix~\ref{sec:repulsive-field} for vector field mathematics and Appendix~\ref{sec:nfz-interaction} for NFZ margin analysis).

\begin{figure}[ht]
\centering
\includegraphics[width=\linewidth]{geofence_diagram}
\caption{Geofence protection visualisation. (a)~Five protection layers: Flight Area hard boundary (emergency RTL), SSSI no-fly zone with 30\,m waypoint buffer, 20\,m speed ramp, 3\,m repulsive field, and hard cutoff. (b)~Repulsive force field that deflects the drone away from the SSSI boundary. (c)~Waypoint correction: the commanded waypoint (red) is shifted to a safe position (green) by the repulsive offset. (d)~Emergency mode switch when the drone enters the SSSI zone.}
\label{fig:geofence_diagram}
\end{figure}

% ═══════════════════════════════════════════════════════════════
\subsection{Computer Vision Pipeline}\label{sec:cv}

The vision subsystem is encapsulated in \texttt{vision.py}, exposing a single interface: \texttt{detect\_in\_image(frame)} $\to$ \texttt{(found, cx, cy, confidence)}. The pipeline proceeds as follows:

\paragraph{Preprocessing.} Lens undistortion (if calibration data are available), colour-space handling (the IMX296 outputs BGR despite its RGB888 label---no conversion applied), resize from $1456\times1088$ to $640\times640$ via bilinear interpolation, and float32 normalisation ($[0,255] \to [0,1]$), yielding input tensor $[1, 640, 640, 3]$.

\paragraph{Inference.} The TFLite model produces output tensor $[1, 5, 8400]$ (8400 anchor-free candidates across three feature pyramid scales). Post-processing applies confidence filtering (threshold 0.2), best-detection selection (greedy single-target), and coordinate rescaling to original frame dimensions. The Ultralytics backend handles non-maximum suppression (NMS) internally; the TFLite backend implements equivalent filtering in \texttt{vision.py}.

\paragraph{Three backends.} The system supports Ultralytics~\cite{yolov8} (development laptop) and TFLite~\cite{tflite} (production Pi), with an NCNN~\cite{ncnn} backend implemented but not yet benchmarked on the Pi. Backend selection is automatic at import time. Both validated backends consume the same model weights and produce identical detection semantics (Appendix~\ref{app:cv-extended} provides full benchmark results and the dual-backend architecture details).

\begin{figure}[ht]
\centering
\includegraphics[width=0.85\textwidth]{cv_pipeline}
\caption{Computer vision pipeline: frame capture, lens undistortion, resize to $640\times640$, TFLite inference, confidence filtering, and coordinate rescaling. The entire pipeline executes in 208\,ms on the Pi~5.}
\label{fig:cv_pipeline}
\end{figure}

\paragraph{Smart detection mode.} An optional \texttt{-{}-smart-detect} flag requires $k$ consecutive positive frames before triggering investigation (default $k{=}3$), suppressing transient false positives at the cost of a 0.6\,s delay.

\paragraph{Performance.} On the Pi~5 with XNNPACK CPU delegate: 206.5\,ms per frame (4.8\,FPS), 100\% detection rate (50/50 bench frames), 0.966 mean confidence (Appendix~\ref{app:cv:benchmarks}). The model was retrained three times---at $640$, $1280$, and $1088$\,px---with the final v2 model achieving mAP$_{50}$ = 0.995 on a dataset of 300 synthetic, 16 real, and 50 negative images (Appendix~\ref{app:cv:training}). At the nominal search speed of 8\,m/s at \SI{35}{\metre} altitude, adjacent frames overlap by approximately 95\% ($1 - 1.67/32.2$), providing 11--17 detection opportunities per target during a single flyover (speed--detection trade-off analysis in Appendix~\ref{app:cv:speed}).

% ═══════════════════════════════════════════════════════════════
\subsection{Search Pattern Generation}\label{sec:search-pattern}

The boustrophedon (lawnmower) pattern~\cite{choset2001coverage} is implemented via a rotate--rasterise--zigzag--rotate-back algorithm in \texttt{planning.py}: the GPS polygon is rasterised into a binary mask, rotated to align with the longest edge (minimising U-turns), stripped at camera-footprint spacing, and inverse-rotated back to GPS coordinates. B\'{e}zier-smoothed U-turns maintain momentum.

\paragraph{Coverage guarantee.}
Strip spacing derives from the thin-lens footprint: $W_g = w_s \cdot h / f$, lane spacing $L = W_g(1 - \alpha)$. At \SI{35}{\metre}: $W_g \approx \SI{32.2}{\metre}$, $L \approx \SI{25.7}{\metre}$ (20\% overlap), guaranteeing every ground point falls within at least one swath under \SIrange{2}{3}{\metre} GPS drift. A 216-configuration sweep confirmed the optimal scan angle at \SIrange{65}{75}{\degree} (Appendix~\ref{sec:path-opt}). Figure~\ref{fig:coverage_vs_time} shows near-linear coverage progression.

\begin{figure}[ht]
  \centering
  \includegraphics[width=0.8\linewidth]{coverage_vs_time}
  \caption{Area coverage progression over time for the boustrophedon search pattern at \SI{35}{\metre} altitude.}
  \label{fig:coverage_vs_time}
\end{figure}

% ═══════════════════════════════════════════════════════════════
\subsection{Mission State Machine}\label{sec:state-machine}

The 20-state FSM (Figure~\ref{fig:state_machine}, Appendix~\ref{app:states}) orchestrates the mission through four phases:

\begin{figure}[htbp]
\centering
\includegraphics[width=0.95\textwidth]{state_machine}
\caption{State transition diagram. Primary mission path flows left to right; engagement branches downward from SEARCH. Manual override (dashed) is accessible from any active state.}
\label{fig:state_machine}
\end{figure}

\begin{itemize}
    \item \textbf{Startup:} INIT $\to$ CONNECTING $\to$ ARMING (GPS fix $\geq$3D, $\geq$6 sats) $\to$ TAKEOFF.
    \item \textbf{Search:} PRE\_WAYPOINTS $\to$ TRANSIT\_TO\_SEARCH $\to$ SEARCH (lawnmower pattern with CV detection; PLB beacon redirect if timer expires).
    \item \textbf{Engagement:} CENTERING (fly to GPS estimate) $\to$ DESCENDING (reduce altitude for close-up view) $\to$ VERIFY (operator Y/N/I with 120\,s timeout; auto-reject resumes search) $\to$ APPROACH (7.5\,m offset) $\to$ HOVER\_TARGET (servo payload release). If the operator rejects or the timeout expires, RETURN\_TO\_SEARCH climbs back to search altitude and resumes the lawnmower pattern from the next unvisited waypoint.
    \item \textbf{Recovery:} RETURN\_TRANSIT $\to$ RETURN\_HOME $\to$ LANDING (triple touchdown detection) $\to$ DONE.
    \item \textbf{Override:} MANUAL mode is reachable from any active state via the M~key, RC switch, or geofence hard-cutoff. All autonomous commands cease; the pilot assumes full control. Re-entering autonomous mode requires an explicit resume command.
\end{itemize}

A detection queue with spatial deduplication (5\,m reject radius), bounded capacity (20 entries), and lazy validation manages multiple targets. Six independent safety mechanisms provide defence-in-depth: manual override, RC kill switch, RTL on link loss (\SI{3}{\second} heartbeat timeout), Flight Area boundary enforcement (emergency RTL), four-layer SSSI geofence, and per-state timeouts. Appendix~\ref{app:states} lists all 20~states with their entry conditions, exit transitions, and timeout values.

% ═══════════════════════════════════════════════════════════════
\subsection{Target Localisation}\label{sec:localisation}

Each detection is converted from pixel coordinates to GPS using the pinhole camera model. The pixel offset $(\Delta u, \Delta v)$ from the image centre is projected to a body-frame displacement via the ground sampling distance:
\begin{equation}\label{eq:gsd-sysdesc}
  \mathrm{GSD} = \frac{w_s \cdot h}{f \cdot W},\qquad
  d_x = \Delta u \cdot \mathrm{GSD},\quad
  d_y = \Delta v \cdot \mathrm{GSD}
\end{equation}
where $w_s = \SI{5.02}{\milli\metre}$ is the sensor width, $f = \SI{5.46}{\milli\metre}$ the calibrated focal length, $h$ the barometric altitude, and $W = 1456$\,px the image width. The body-frame displacement $(d_x, d_y)$ is rotated into North--East coordinates using the drone's heading and converted to latitude/longitude on the WGS\,84 ellipsoid. Appendix~\ref{sec:gps-deep} documents the complete 15-stage detection-to-GPS pipeline with a full flowchart, error budget decomposition by flight phase, and roll/pitch geometry analysis (Figure~\ref{fig:pipeline-flow}). Multiple observations are fused through:

\begin{itemize}
    \item \textbf{Spatial clustering} (30\,m merge radius) for multi-target tracking;
    \item \textbf{Inverse-variance weighting} --- altitude factor $w_\mathrm{alt} = (h_\mathrm{ref}/h)^2$ and centrality factor $w_\mathrm{cen} = 1 + 4(1 - d/d_\mathrm{max})$;
    \item \textbf{Kalman filter} on the $[\phi, \lambda]$ state vector with measurement noise scaled by inverse observation weight, converging within 5--10 observations.
\end{itemize}

Video replay analysis (DJI flight footage, $1456 \times 1088$ at 30\,fps) yielded CEP$_{50}$ = 2.3\,m and a maximum error of 16.5\,m (high-speed outlier). The characteristic directional elongation stems from GPS receiver latency (100--200\,ms), which multi-pass averaging in the lawnmower pattern partially mitigates. Six independent error sources contribute to a single-observation RSS uncertainty of \SI{2.9}{\metre} at \SI{35}{\metre} altitude, dominated by GPS receiver noise (\SI{2.3}{\metre}) and timing lag (\SI{1.5}{\metre}). After inverse-variance fusion of 15+ multi-altitude observations, the fused estimate converges to sub-metre accuracy, exceeding the 5--10\,m landing requirement (R07) with $>$99\% probability. Appendix~\ref{sec:gps-deep} provides the full error budget, convergence plots, and fusion strategy comparison.

\begin{figure}[ht]
\centering
\begin{subfigure}[t]{0.48\textwidth}
  \centering
  \includegraphics[width=\linewidth]{gps_bullseye}
  \caption{GPS estimation error scatter (bullseye). CEP$_{50}$ = 2.3\,m shown as dashed circle.}
  \label{fig:gps_bullseye}
\end{subfigure}
\hfill
\begin{subfigure}[t]{0.48\textwidth}
  \centering
  \includegraphics[width=\linewidth]{gps_error_direction}
  \caption{Directional error distribution showing heading-aligned elongation from GPS timing lag.}
  \label{fig:gps_error_direction}
\end{subfigure}
\caption{Target localisation accuracy from DJI flight video replay analysis.}
\label{fig:gps_accuracy}
\end{figure}

% ═══════════════════════════════════════════════════════════════
\subsection{Ground Station}\label{sec:groundstation}

The ground station is a browser-based dashboard served by a threaded HTTP server on the Pi at port~8090, accessible from any device on the local Wi-Fi network. It provides three panels: (1)~live MJPEG video stream with detection overlays; (2)~2D GPS grid showing detection clusters, drone position, and search path; and (3)~command buttons (investigate, confirm, reject, land, RTL, manual override). MJPEG was chosen over H.264/WebRTC for its minimal client complexity (single \texttt{<img>} tag), zero transcoding load, and per-frame decodability. Measured latency is \SI{300}{\milli\second} at \SI{0.3}{\mega\byte\per\second}. The Pi auto-detects headless operation (no \texttt{DISPLAY}), enabling SSH deployment without X11 forwarding (Appendix~\ref{sec:streaming-arch}).

% ═══════════════════════════════════════════════════════════════
\subsection{Simulation Framework}\label{sec:simulation}

Four simulation levels provide progressive fidelity (Appendix~\ref{sec:sim-validation}): (1)~interactive simulator (keyboard flight over satellite map with ground-truth comparison); (2)~SITL integration (full ArduPilot autopilot with simulated camera view); (3)~SITL with real webcam; (4)~DJI video replay with synchronised SRT telemetry. All four levels share the same detection model, GPS estimation mathematics, and MAVLink command interface, so transitioning to real hardware requires only configuration changes. Simulation alone caught 12 defects---including a MAVLink race condition during takeoff, incorrect latitude scaling, auto-disarm timeout, and GPS timing lag---all resolved before any hardware integration.

\begin{figure}[ht]
\centering
\includegraphics[width=0.9\textwidth]{mission_timeline}
\caption{Mission timeline showing state transitions during a simulated end-to-end SAR mission. The horizontal axis represents elapsed time; coloured blocks indicate active states.}
\label{fig:mission_timeline}
\end{figure}

